% =============================================================================
% 3. PROBLEM FORMULATION AND ENGINEERING PRELIMINARIES
% =============================================================================
\section{Problem Formulation and Engineering Preliminaries}
\label{sec:problem_formulation}

\subsection{Sequential Edge Few-Shot Task Formulation}
\label{sec:task_formulation}
Field diagnosis requires identifying crop diseases directly on mobile devices as leaves are inspected in the field. When deployed to a new crop or region, the system must calibrate using only a small reference set of labeled leaves collected on site:
\begin{equation}
    \mathcal{S} = \{(\mathbf{x}_i, y_i)\}_{i=1}^{N_s}, \qquad N_s = C \times K
\end{equation}
where $C$ is the number of target pathology classes (e.g., $C \in \{5, 11, 27\}$) and $K$ is the number of labeled examples per class (typically $K = 5$). These local diseases are unseen during pre-training ($\mathcal{Y}_{\text{base}} \cap \mathcal{Y}_{\text{novel}} = \emptyset$).

During active scouting, leaves are observed sequentially as an online stream:
\begin{equation}
    \mathcal{Q}_{\text{stream}} = \{\mathbf{x}_t\}_{t=1}^{T}, \qquad \mathbf{x}_t \in \mathbb{R}^{H \times W \times 3}
\end{equation}
The system must output an immediate diagnosis $\hat{y}_t \in \{1, \dots, C\}$ for each incoming frame $\mathbf{x}_t$, without buffering queries or waiting for future images.

Performing this diagnosis on edge hardware faces strict computational limits. Fine-tuning deep backbones on device requires extra memory and causes processor slowdown. The model must therefore rely on frozen foundation encoders ($\nabla_\theta = \emptyset$). We extract complementary visual and biological features using DINOv2-Small ($d_1 = 384$) and BioCLIP ($d_2 = 512$), yielding a fixed joint descriptor $\mathbf{z} \in \mathbb{R}^d$ ($d = 896$) without gradient calculations during field adaptation.

\subsection{Engineering and Operational Constraints}
\label{sec:engineering_constraints}
Because remote crop fields often lack internet connectivity, foliar diagnosis must run entirely on edge devices rather than on cloud servers \cite{shoaib2023advanced}. We define the on-device diagnostic objective as a constrained optimization problem that maximizes accuracy subject to four operational engineering boundaries:
\begin{equation}
\begin{aligned}
    \max_{\mathbf{W}, \mathbf{b}, \, \tau_{\text{conf}}} \quad & \mathbb{E}_{(\mathbf{x}_t, y_t)}\big[\mathbb{I}(\hat{y}_t = y_t)\big] \\
    \text{subject to} \quad & \begin{cases}
        \text{FNR}(\hat{y}) \le \epsilon_{\text{esc}} & \text{(agronomic safety)} \\
        \overline{L}_{\text{inf}}(\tau_{\text{conf}}) \le L_{\max} & \text{(real-time stream latency)} \\
        T_{\text{adapt}}(\mathcal{S}) \le T_{\max}, \; \nabla_\theta = \emptyset & \text{(gradient-free calibration)} \\
        \text{Mem}_{\text{PSS}} \le M_{\max} & \text{(mobile thermal/RAM budget)}
    \end{cases}
\end{aligned}
\label{eq:constrained_objective}
\end{equation}
where $\mathbf{W} \in \mathbb{R}^{C \times d}$ and $\mathbf{b} \in \mathbb{R}^C$ denote the analytical projection parameters, and $\tau_{\text{conf}}$ governs the two-tier escalation gate.

While the exact operational parameters $(\epsilon_{\text{esc}}, L_{\max}, T_{\max}, M_{\max})$ depend on the scouting platform and target crop pathology, we establish representative target boundaries for our edge deployment:
\begin{enumerate}
    \item \textbf{Agronomic Disease Escape Limit ($\text{FNR} \le \epsilon_{\text{esc}}$):} In plant protection, misclassifying an infected leaf as healthy allows pathogens to spread across the field unnoticed. Under integrated pest management principles \cite{richard2022control}, preventing secondary infection outbreaks requires minimizing missed infections. We therefore target a false-negative rate ($\text{FNR}$) below $\epsilon_{\text{esc}} = 1.0\%$.
    \item \textbf{Real-Time Stream Latency ($\overline{L}_{\text{inf}} \le L_{\max}$):} Field scouting platforms require continuous scanning to inspect crop canopies. Maintaining an inspection rate of $\ge 6$ frames per second requires mean per-frame latency to remain within $L_{\max} = 170$\,ms to prevent buffer overflow and dropped frames \cite{sharif2026real}.
    \item \textbf{Sub-Second On-Device Calibration ($T_{\text{adapt}} \le T_{\max}$):} When encountering a new regional pathogen variant in the field, calibrating the model with a few reference leaves ($K=5$) must finish within sub-second time ($T_{\max} \le 1.0$\,s) so scouting can resume immediately. In addition, backpropagation and optimizer tracking are disallowed on device ($\nabla_\theta = \emptyset$) to prevent memory bloat and processor slowdown.
    \item \textbf{Memory Footprint ($\text{Mem}_{\text{PSS}} \le M_{\max}$):} Mobile edge terminals share limited random-access memory (RAM) across the operating system and camera drivers. Setting a budget of $M_{\max} = 500$\,MB peak Proportional Set Size (PSS) ensures the diagnostic engine avoids memory-pressure termination and operates safely on fanless embedded hardware.
\end{enumerate}

\subsection{Statistical Preliminaries: Small-Sample Covariance Degeneracy}
\label{sec:covariance_degeneracy}
In foliar image representations, feature variance is heavily non-uniform across latent dimensions. Leaf imagery naturally shares dominant morphological structures---such as healthy green chlorophyll pigmentation, background venation patterns, and specular surface reflections---that span high-variance directions in the feature space. When classification relies on isotropic metrics such as unweighted Euclidean distance or cosine similarity:
\begin{equation}
    \hat{y}_t = \arg\max_{c \in \{1, \dots, C\}} \frac{\mathbf{z}_t^T \boldsymbol{\mu}_c}{\|\mathbf{z}_t\|_2 \|\boldsymbol{\mu}_c\|_2}, \qquad \boldsymbol{\mu}_c = \frac{1}{|\mathcal{S}_c|}\sum_{\mathbf{x} \in \mathcal{S}_c} \mathbf{z}(\mathbf{x})
    \label{eq:isotropic_cosine}
\end{equation}
all feature coordinates are weighted equally ($\mathbf{\Sigma} = \sigma^2 \mathbf{I}$). Consequently, distance calculations are dominated by shared vegetative background variance rather than the subtle, localized lesion characteristics critical for disease diagnosis.

Fisher Linear Discriminant Analysis resolves this issue by projecting features through the inverse within-class covariance $\hat{\mathbf{\Sigma}}^{-1}$, which suppresses high-variance shared background directions and amplifies discriminative lesion signals. Computing this projection requires the empirical pooled within-class scatter matrix:
\begin{equation}
    \mathbf{S}_W = \frac{1}{N - C} \sum_{c=1}^C \sum_{\mathbf{z} \in \mathcal{S}_c} (\mathbf{z} - \boldsymbol{\mu}_c)(\mathbf{z} - \boldsymbol{\mu}_c)^T \in \mathbb{R}^{d \times d}
    \label{eq:empirical_scatter}
\end{equation}
However, in realistic few-shot deployment, the total support sample size $N$ is substantially smaller than the latent dimension $d$ ($N \ll d$). Across our target evaluation settings:
\begin{itemize}
    \item In 5-way 5-shot field diagnosis, $N = 25$ versus $d = 896$, yielding $\text{rank}(\mathbf{S}_W) \le 20$.
    \item On the 11-class PlantVillage benchmark without augmentation, $N = 55$ versus $d = 896$, yielding $\text{rank}(\mathbf{S}_W) \le 44$.
    \item Even after multi-view support augmentation ($N = 440$), $\text{rank}(\mathbf{S}_W) \le 429 < 896$.
\end{itemize}
Because $\text{rank}(\mathbf{S}_W) \le N - C \ll d$, the empirical scatter matrix $\mathbf{S}_W$ has a nullity of at least $d - (N - C) > 0$, making it rank-deficient, strictly non-invertible ($\det(\mathbf{S}_W) = 0$), and prone to severe numerical instability under naive diagonal loading. This fundamental singularity motivates the analytical shrinkage and uncertainty-gated architecture detailed in Section~\ref{sec:methodology}. Table~\ref{tab:notation} summarizes the mathematical notations used throughout this paper.

\begin{table*}[htbp]
\caption{Summary of mathematical notations and descriptions used in BioDINO-Edge.}
\label{tab:notation}
\centering
\footnotesize
\begin{tabularx}{\textwidth}{lX}
\toprule
\textbf{Notation} & \textbf{Mathematical Description and Operational Meaning} \\
\midrule
$C, K, N_s$ & Number of pathology classes, shots per class ($K=5$), and total support size ($N_s = C \cdot K$). \\
$\phi_{\text{dino}}(\cdot), \phi_{\text{bio}}(\cdot)$ & Frozen INT8 DINOv2-Small ($d_1 = 384$) and BioCLIP ViT-B/16 ($d_2 = 512$) encoders. \\
$\mathbf{z} \in \mathbb{R}^{896}$ & Concatenated joint foliar descriptor: $\mathbf{z} = [\mathbf{f}_{\text{bio}} \,\|\, \mathbf{f}_{\text{dino}}] \in \mathbb{R}^{896}$. \\
$\mathbf{P}_{\text{dino}}, \mathbf{P}_{\text{text}}$ & Tier-1 visual prototypes ($\mathbb{R}^{C \times 384}$) and cached botanical text prototypes ($\mathbb{R}^{C \times 512}$). \\
$\mathbf{S}_W, \hat{\mathbf{\Sigma}}$ & Empirical pooled within-class scatter matrix and Ledoit--Wolf regularized covariance matrix. \\
$\rho^*, \nu$ & Analytical shrinkage intensity factor $\rho^* \in (0, 1]$ and mean trace shrinkage target $\nu = \frac{1}{d}\text{Tr}(\mathbf{S}_W)$. \\
$\mathbf{W}, \mathbf{b}$ & Fisher-LDA discriminant projection matrix ($\mathbb{R}^{C \times 896}$) and class bias vector ($\mathbb{R}^C$). \\
$\mathcal{C}(\mathbf{x}), \tau_{\text{conf}}$ & Dual-metric uncertainty gate (Shannon entropy and top-2 margin) and escalation threshold ($\tau_{\text{conf}} = 0.35$). \\
$w_{\text{lda}}$ & Multimodal consensus blending weight ($w_{\text{lda}} = 0.75$, weighting visual Fisher against text posteriors). \\
$\gamma$ & Inverse temperature logit scaling parameter ($\gamma = 20.0$). \\
$\mathcal{Q}_{\text{anchor}}^{(c)}$ & High-confidence streaming query buffer ($\max_c p_c \ge 0.70$) for background transductive recalibration. \\
\bottomrule
\end{tabularx}
\end{table*}
